\section*{Background}

Hurricane Helene caused flooding and wind damage across parts of the southeastern United States in late September 2024. Such a shock can force firms to balance long-term investment against immediate funding needs. Research and development (R\&D) is particularly relevant because its benefits often arrive well after the initial spending.

Dessaint and Matray (2017) find that firms near a hurricane disaster area, but not directly affected, temporarily increase cash holdings and discuss hurricane risk more often in regulatory filings. They interpret this as salience: a nearby disaster makes liquidity risk more prominent in managers' decisions. Their study motivates examining investment alongside cash, although directly exposed firms also face damage and operational disruption.

The research question is: how does R\&D investment change after Helene among firms headquartered in affected counties, relative to firms elsewhere, and is a decline accompanied by higher cash holdings? The joint pattern would be consistent with greater emphasis on liquidity, but would not by itself establish a precautionary motive.
